\documentclass[11pt,oneside]{book}
\usepackage[letterpaper,margin=1.1in]{geometry}
\usepackage{fontspec}
\usepackage{microtype}
\usepackage{xcolor}
\usepackage{enumitem}
\usepackage{titlesec}
\usepackage{hyperref}
\setmainfont{Latin Modern Roman}
\definecolor{linkblue}{HTML}{164E8A}
\hypersetup{colorlinks=true,linkcolor=linkblue,urlcolor=linkblue,citecolor=linkblue,pdfauthor={We The Nameless},pdftitle={Dude Index: The Dudetheyreonamystic History}}
\setlist[description]{style=nextline,leftmargin=0pt,labelsep=0pt,font=\normalfont\bfseries,itemsep=0.45\baselineskip,parsep=0pt}
\titleformat{\chapter}[display]{\normalfont\huge\bfseries}{\chaptertitlename\ \thechapter}{0.5em}{}
\titleformat{\section}{\normalfont\Large\bfseries}{\thesection}{0.75em}{}
\titlespacing*{\section}{0pt}{1.4\baselineskip}{0.55\baselineskip}
\newcommand{\chronologysection}[3]{\section[#1: #2]{#1\\\large #2\\\normalsize\normalfont #3}}

\begin{document}
\frontmatter
\begin{titlepage}
\centering
\vspace*{0.18\textheight}
{\Huge\bfseries Dude Index\par}
\vspace{1em}
{\LARGE The Dudetheyreonamystic History\par}
\vspace{2em}
{\large A plan for Dudetheyreontome and the Dudetheyreonamystic History\par}
\vfill
{\large Generated from \texttt{Dude History.md}\par}
\end{titlepage}

\chapter*{Plan}
\addcontentsline{toc}{chapter}{Plan}
\begin{description}
  \item[Subject] The life of the Bible after the end of II Kings.
  \item[Method] Assemble short extracts from scriptures, apocrypha, pseudepigrapha,
  archives, translations, commentaries, polemics, archaeological discoveries,
  and modern scholarship—exactly as the Deuteronomistic Historian assembled
  older sources into a new history.
  \item[Central plot] The old national literature survives the destruction of its
  kingdom; becomes Torah and Bible; divides into variant textual families and
  rival canons; generates Judaism, Christianity, and Islam; is stabilized by
  scribes; reopened by philology; divided into sources by criticism; and placed
  back into its ancient material world by archaeology.
\end{description}

\tableofcontents
\mainmatter
\chapter{Chronology}

\chronologysection{587/586 BCE}{II Kings 25}{Jerusalem}
\begin{description}
  \item[Date described] 587/586 BCE
  \item[Place / stage] Jerusalem
  \item[Best source or object] \textbf{II Kings 25}
  \item[Suggested cut] 25:8–12, 18–21, 27–30
  \item[What happens to the Bible] The old history ends: temple burned, monarchy destroyed, people deported—but Jehoiachin survives in Babylon, leaving the story barely open.
  \item[Approximate date written / made] Sixth century BCE, with later editing
\end{description}

\chronologysection{597–570 BCE}{Ezekiel}{Babylonia}
\begin{description}
  \item[Date described] 597–570 BCE
  \item[Place / stage] Babylonia
  \item[Best source or object] \textbf{Ezekiel}
  \item[Suggested cut] 1:1–3; 8:1–12; 11:14–21; 18:1–4; 37:1–14
  \item[What happens to the Bible] Prophecy detaches YHWH from the land and temple. Israel’s god can travel, speak, judge, and restore in exile.
  \item[Approximate date written / made] Mostly sixth century BCE
\end{description}

\chronologysection{c. 594 BCE}{Jeremiah’s letter to the exiles}{Jerusalem → Babylon}
\begin{description}
  \item[Date described] c. 594 BCE
  \item[Place / stage] Jerusalem → Babylon
  \item[Best source or object] \textbf{Jeremiah’s letter to the exiles}
  \item[Suggested cut] Jeremiah 29:1–14
  \item[What happens to the Bible] Exile becomes a durable way of life: build houses, plant gardens, marry, multiply, and seek the city’s welfare.
  \item[Approximate date written / made] Sixth century BCE, edited later
\end{description}

\chronologysection{586 BCE}{Lamentations}{Jerusalem}
\begin{description}
  \item[Date described] 586 BCE
  \item[Place / stage] Jerusalem
  \item[Best source or object] \textbf{Lamentations}
  \item[Suggested cut] 1:1–4; 2:10–15; 5:19–22
  \item[What happens to the Bible] National defeat becomes portable liturgy. The ruined city survives as a book that can be mourned anywhere.
  \item[Approximate date written / made] Sixth century BCE
\end{description}

\chronologysection{After 586 BCE}{Psalm 137}{Babylon}
\begin{description}
  \item[Date described] After 586 BCE
  \item[Place / stage] Babylon
  \item[Best source or object] \textbf{Psalm 137}
  \item[Suggested cut] Entire psalm
  \item[What happens to the Bible] The first great diaspora poem: memory of Zion, refusal to assimilate completely, grief turning into revenge fantasy.
  \item[Approximate date written / made] Sixth–fifth century BCE
\end{description}

\chronologysection{586–539 BCE}{Isaiah 40–55}{Babylon}
\begin{description}
  \item[Date described] 586–539 BCE
  \item[Place / stage] Babylon
  \item[Best source or object] \textbf{Isaiah 40–55}
  \item[Suggested cut] 40:1–8; 44:24–45:7; 55:10–11
  \item[What happens to the Bible] Israel’s literature reinterprets catastrophe: Cyrus becomes YHWH’s anointed, and the divine word—not the monarchy—becomes permanent.
  \item[Approximate date written / made] Sixth century BCE
\end{description}

\chronologysection{Sixth–fifth centuries BCE}{Al-Yahudu tablets}{Babylonia}
\begin{description}
  \item[Date described] Sixth–fifth centuries BCE
  \item[Place / stage] Babylonia
  \item[Best source or object] \textbf{Al-Yahudu tablets}
  \item[Suggested cut] A few contracts naming Judean families, fields, debts, and inherited YHWH-names
  \item[What happens to the Bible] The exiles become farmers, tenants, merchants, and imperial subjects. Biblical exile meets ordinary paperwork.
  \item[Approximate date written / made] c. 572–477 BCE
\end{description}

\chronologysection{539–515 BCE}{Ezra 1–6}{Persia and Yehud}
\begin{description}
  \item[Date described] 539–515 BCE
  \item[Place / stage] Persia and Yehud
  \item[Best source or object] \textbf{Ezra 1–6}
  \item[Suggested cut] 1:1–8; 3:10–13; 4:7–16; 6:13–18
  \item[What happens to the Bible] Imperial decrees, archives, and letters become part of sacred history. The rebuilt temple replaces the lost kingdom.
  \item[Approximate date written / made] Persian-period materials; edited fourth–third centuries BCE
\end{description}

\chronologysection{520–518 BCE}{Haggai and Zechariah}{Jerusalem}
\begin{description}
  \item[Date described] 520–518 BCE
  \item[Place / stage] Jerusalem
  \item[Best source or object] \textbf{Haggai and Zechariah}
  \item[Suggested cut] Haggai 1:2–11; Zechariah 4:6–10
  \item[What happens to the Bible] Prophets turn rebuilding into revelation. The temple is restored, but the Davidic monarchy is not.
  \item[Approximate date written / made] Late sixth century BCE
\end{description}

\chronologysection{Fifth century BCE}{Priestly Torah}{Babylonia / Judah}
\begin{description}
  \item[Date described] Fifth century BCE
  \item[Place / stage] Babylonia / Judah
  \item[Best source or object] \textbf{Priestly Torah}
  \item[Suggested cut] Genesis 1:1–2:4a; Genesis 17; Exodus 25:1–9; Leviticus 16:29–34
  \item[What happens to the Bible] Creation, covenant, purity, sacrifice, calendar, and portable sanctuary are organized into an immense priestly system.
  \item[Approximate date written / made] Commonly exilic or Persian period; precise model disputed
\end{description}

\chronologysection{c. 460–400 BCE}{Elephantine archives}{Elephantine, Egypt}
\begin{description}
  \item[Date described] c. 460–400 BCE
  \item[Place / stage] Elephantine, Egypt
  \item[Best source or object] \textbf{Elephantine archives}
  \item[Suggested cut] One marriage contract; one property document; one letter mentioning YHW
  \item[What happens to the Bible] An actual Judean diaspora has its own temple, priests, mixed households, Aramaic legal culture, and limited dependence on Jerusalem.
  \item[Approximate date written / made] Fifth century BCE
\end{description}

\chronologysection{419/418 BCE}{Elephantine Passover Letter}{Elephantine}
\begin{description}
  \item[Date described] 419/418 BCE
  \item[Place / stage] Elephantine
  \item[Best source or object] \textbf{Elephantine Passover Letter}
  \item[Suggested cut] Surviving lines concerning unleavened bread, dates, and prohibited work
  \item[What happens to the Bible] Festival law travels through imperial correspondence. Torah-like practice is negotiated before the Pentateuch is visibly fixed everywhere.
  \item[Approximate date written / made] 419/418 BCE
\end{description}

\chronologysection{410–407 BCE}{Temple Petition}{Elephantine}
\begin{description}
  \item[Date described] 410–407 BCE
  \item[Place / stage] Elephantine
  \item[Best source or object] \textbf{Temple Petition}
  \item[Suggested cut] Description of the destruction of YHW’s temple and request to rebuild
  \item[What happens to the Bible] Elephantine appeals to authorities in both Jerusalem and Samaria. There is not yet one uncontested center speaking for all Judeans.
  \item[Approximate date written / made] 407 BCE
\end{description}

\chronologysection{Fifth–fourth centuries BCE}{Ezra 7–10; Nehemiah 8–10}{Persia / Jerusalem}
\begin{description}
  \item[Date described] Fifth–fourth centuries BCE
  \item[Place / stage] Persia / Jerusalem
  \item[Best source or object] \textbf{Ezra 7–10; Nehemiah 8–10}
  \item[Suggested cut] Ezra 7:6–10; Nehemiah 8:1–12; 9:1–3
  \item[What happens to the Bible] A written Torah is publicly read, explained, and made the constitutional center of the restored community.
  \item[Approximate date written / made] Fifth-century memoir traditions; edited later
\end{description}

\chronologysection{Legendary restoration}{2 Esdras 14 / 4 Ezra}{Babylon / Jerusalem}
\begin{description}
  \item[Date described] Legendary restoration
  \item[Place / stage] Babylon / Jerusalem
  \item[Best source or object] \textbf{2 Esdras 14 / 4 Ezra}
  \item[Suggested cut] 14:19–48
  \item[What happens to the Bible] Ezra dictates the destroyed scriptures anew: twenty-four public books and seventy secret books. A myth of total scriptural reconstruction.
  \item[Approximate date written / made] c. 90–100 CE, describing the age of Ezra
\end{description}

\chronologysection{Fourth–third centuries BCE}{Chronicles}{Judah}
\begin{description}
  \item[Date described] Fourth–third centuries BCE
  \item[Place / stage] Judah
  \item[Best source or object] \textbf{Chronicles}
  \item[Suggested cut] 1 Chronicles 9:1–2; 2 Chronicles 36:22–23
  \item[What happens to the Bible] Samuel–Kings is rewritten: David becomes temple founder, Levites dominate, northern material shrinks, and the ending points toward restoration.
  \item[Approximate date written / made] Fourth–third centuries BCE
\end{description}

\chronologysection{Third century BCE}{Septuagint Pentateuch}{Alexandria}
\begin{description}
  \item[Date described] Third century BCE
  \item[Place / stage] Alexandria
  \item[Best source or object] \textbf{Septuagint Pentateuch}
  \item[Suggested cut] Genesis 1; Exodus 3; Deuteronomy 32 in Hebrew and Greek side by side
  \item[What happens to the Bible] Torah enters Greek. Translation creates new meanings, new textual forms, and eventually a second major scriptural civilization.
  \item[Approximate date written / made] Pentateuch translated mainly third century BCE
\end{description}

\chronologysection{Imagined third century BCE}{Letter of Aristeas}{Alexandria / Jerusalem}
\begin{description}
  \item[Date described] Imagined third century BCE
  \item[Place / stage] Alexandria / Jerusalem
  \item[Best source or object] \textbf{Letter of Aristeas}
  \item[Suggested cut] §§9–11; 301–311 or another brief account of the translators
  \item[What happens to the Bible] The Greek Torah receives a charter myth: royal library, seventy-two translators, miraculous agreement, and an authoritative translation.
  \item[Approximate date written / made] Probably second century BCE
\end{description}

\chronologysection{Third–second centuries BCE}{Tobit}{Babylonia / Media}
\begin{description}
  \item[Date described] Third–second centuries BCE
  \item[Place / stage] Babylonia / Media
  \item[Best source or object] \textbf{Tobit}
  \item[Suggested cut] 1:1–22; 4:5–19; 6:1–9; 14:3–7
  \item[What happens to the Bible] Exile becomes family romance: burial, almsgiving, endogamy, angels, demons, medicine, travel, and hope for Jerusalem.
  \item[Approximate date written / made] c. 225–175 BCE
\end{description}

\chronologysection{Second century BCE}{Esther}{Persia}
\begin{description}
  \item[Date described] Second century BCE
  \item[Place / stage] Persia
  \item[Best source or object] \textbf{Esther}
  \item[Suggested cut] Esther 3:8–15; 4:13–17; 9:20–28
  \item[What happens to the Bible] A diaspora community survives through court access, counterviolence, memory, and a new festival—with God almost entirely offstage.
  \item[Approximate date written / made] Probably fourth–third century BCE; dating disputed
\end{description}

\chronologysection{Second–first centuries BCE}{Greek Additions to Esther}{Persia}
\begin{description}
  \item[Date described] Second–first centuries BCE
  \item[Place / stage] Persia
  \item[Best source or object] \textbf{Greek Additions to Esther}
  \item[Suggested cut] Esther’s prayer and Mordecai’s dream
  \item[What happens to the Bible] Greek editors add prayers, dreams, divine providence, and official decrees, making the strangely secular Hebrew story more conventionally biblical.
  \item[Approximate date written / made] Second–first centuries BCE
\end{description}

\chronologysection{Second century BCE}{Artapanus, Demetrius, Ezekiel the Tragedian}{Egypt}
\begin{description}
  \item[Date described] Second century BCE
  \item[Place / stage] Egypt
  \item[Best source or object] \textbf{Artapanus, Demetrius, Ezekiel the Tragedian}
  \item[Suggested cut] One surviving fragment from each
  \item[What happens to the Bible] Greek-speaking Jews rewrite biblical history as chronology, philosophy, ethnography, and stage tragedy. Moses becomes a world-civilizing hero.
  \item[Approximate date written / made] Third–second centuries BCE
\end{description}

\chronologysection{c. 200–175 BCE}{Ben Sira and the translator’s prologue}{Jerusalem / Alexandria}
\begin{description}
  \item[Date described] c. 200–175 BCE
  \item[Place / stage] Jerusalem / Alexandria
  \item[Best source or object] \textbf{Ben Sira and the translator’s prologue}
  \item[Suggested cut] Sirach prologue; 24:1–23; 44:1–15
  \item[What happens to the Bible] “Law, Prophets, and other ancestral books” are translated into Greek. The shape of a larger scriptural collection becomes visible.
  \item[Approximate date written / made] Hebrew c. 200–175 BCE; Greek after 132 BCE
\end{description}

\chronologysection{Third–second centuries BCE}{1 Enoch: Book of the Watchers}{Judea}
\begin{description}
  \item[Date described] Third–second centuries BCE
  \item[Place / stage] Judea
  \item[Best source or object] \textbf{1 Enoch: Book of the Watchers}
  \item[Suggested cut] 1 Enoch 1:1–9; 6–10
  \item[What happens to the Bible] Genesis 6 explodes into an entire cosmic history. A few obscure biblical lines generate angels, giants, forbidden knowledge, and apocalypse.
  \item[Approximate date written / made] Third century BCE, with layers
\end{description}

\chronologysection{Second century BCE}{Jubilees}{Judea}
\begin{description}
  \item[Date described] Second century BCE
  \item[Place / stage] Judea
  \item[Best source or object] \textbf{Jubilees}
  \item[Suggested cut] 1:1–5; 6:17–38; 15:25–34
  \item[What happens to the Bible] Genesis and Exodus are rewritten as a revelation already dictated in heaven, with stricter calendar, purity, circumcision, and ethnic boundaries.
  \item[Approximate date written / made] c. 160–140 BCE
\end{description}

\chronologysection{175–164 BCE}{Daniel 7–12}{Judea}
\begin{description}
  \item[Date described] 175–164 BCE
  \item[Place / stage] Judea
  \item[Best source or object] \textbf{Daniel 7–12}
  \item[Suggested cut] 7:1–14; 9:20–27; 12:1–4
  \item[What happens to the Bible] Current imperial persecution is coded as ancient prediction. Biblical history becomes apocalypse, angelic timetable, and resurrection hope.
  \item[Approximate date written / made] c. 167–164 BCE
\end{description}

\chronologysection{167–160 BCE}{1 Maccabees}{Judea}
\begin{description}
  \item[Date described] 167–160 BCE
  \item[Place / stage] Judea
  \item[Best source or object] \textbf{1 Maccabees}
  \item[Suggested cut] 1:10–15, 41–64; 2:15–28; 4:36–59
  \item[What happens to the Bible] A new biblical-style history explains resistance, martyrdom, temple rededication, and the rise of the Hasmoneans—without becoming canonical in Judaism.
  \item[Approximate date written / made] c. 100 BCE, probably from Hebrew
\end{description}

\chronologysection{180–161 BCE}{2 Maccabees}{Judea / Egypt}
\begin{description}
  \item[Date described] 180–161 BCE
  \item[Place / stage] Judea / Egypt
  \item[Best source or object] \textbf{2 Maccabees}
  \item[Suggested cut] 2:19–32; 6:18–31; 7:1–23; 15:37–39
  \item[What happens to the Bible] History becomes epitome, martyr legend, bodily resurrection, heavenly intervention, and explicit reflection on the editor’s craft.
  \item[Approximate date written / made] Late second century BCE
\end{description}

\chronologysection{Second–first centuries BCE}{Genesis Apocryphon}{Judea}
\begin{description}
  \item[Date described] Second–first centuries BCE
  \item[Place / stage] Judea
  \item[Best source or object] \textbf{Genesis Apocryphon}
  \item[Suggested cut] Columns II–V on Noah’s suspicious birth; XIX–XXII on Abraham
  \item[What happens to the Bible] Biblical figures tell their own stories in first-person Aramaic. Scripture remains open to expansion, correction, and imaginative filling of gaps.
  \item[Approximate date written / made] c. first century BCE, using older traditions
\end{description}

\chronologysection{Second–first centuries BCE}{Community Rule / Damascus Document}{Qumran}
\begin{description}
  \item[Date described] Second–first centuries BCE
  \item[Place / stage] Qumran
  \item[Best source or object] \textbf{Community Rule / Damascus Document}
  \item[Suggested cut] 1QS 1:1–15; CD 1:1–12
  \item[What happens to the Bible] A sect reads prophecy as its own history, builds a covenant community, and treats interpretation as continuing revelation.
  \item[Approximate date written / made] Second–first centuries BCE
\end{description}

\chronologysection{First century BCE}{Habakkuk Pesher}{Qumran}
\begin{description}
  \item[Date described] First century BCE
  \item[Place / stage] Qumran
  \item[Best source or object] \textbf{Habakkuk Pesher}
  \item[Suggested cut] 1QpHab VII–VIII
  \item[What happens to the Bible] “The prophet did not know what he meant.” Scripture’s true referent is the sect’s present, disclosed by the inspired interpreter.
  \item[Approximate date written / made] First century BCE
\end{description}

\chronologysection{c. 50 BCE–50 CE}{Wisdom of Solomon}{Alexandria}
\begin{description}
  \item[Date described] c. 50 BCE–50 CE
  \item[Place / stage] Alexandria
  \item[Best source or object] \textbf{Wisdom of Solomon}
  \item[Suggested cut] 2:10–24; 7:22–30; 11:15–16
  \item[What happens to the Bible] Biblical wisdom fuses with Greek philosophy, immortality, cosmic order, and an allegorical retelling of Exodus.
  \item[Approximate date written / made] Late first century BCE–early first century CE
\end{description}

\chronologysection{c. 50 BCE–100 CE}{Joseph and Aseneth}{Egypt or Syria}
\begin{description}
  \item[Date described] c. 50 BCE–100 CE
  \item[Place / stage] Egypt or Syria
  \item[Best source or object] \textbf{Joseph and Aseneth}
  \item[Suggested cut] Aseneth’s conversion and heavenly meal
  \item[What happens to the Bible] Joseph’s unexplained Egyptian marriage becomes a large romance about conversion, identity, food, and entrance into Israel.
  \item[Approximate date written / made] First century BCE or CE; disputed
\end{description}

\chronologysection{38–41 CE}{Philo, Against Flaccus and Embassy to Gaius}{Alexandria / Rome}
\begin{description}
  \item[Date described] 38–41 CE
  \item[Place / stage] Alexandria / Rome
  \item[Best source or object] \textbf{Philo, Against Flaccus and Embassy to Gaius}
  \item[Suggested cut] Short account of the Alexandrian pogrom; Caligula’s statue crisis
  \item[What happens to the Bible] A Greek philosophical Bible meets imperial politics. Allegorical exegesis becomes a way of inhabiting both Moses and Greek culture.
  \item[Approximate date written / made] c. 40–50 CE
\end{description}

\chronologysection{c. 30 CE}{Mark}{Galilee and Judea}
\begin{description}
  \item[Date described] c. 30 CE
  \item[Place / stage] Galilee and Judea
  \item[Best source or object] \textbf{Mark}
  \item[Suggested cut] Mark 1:1–15; 8:27–33; 13:1–8; 15:33–39
  \item[What happens to the Bible] Israel’s scriptures are recut around Jesus. A failed messianic execution becomes “good news,” and an old prophetic genre mutates into gospel.
  \item[Approximate date written / made] c. 65–75 CE
\end{description}

\chronologysection{c. 30–60 CE}{Paul: minimal anthology}{Eastern Mediterranean}
\begin{description}
  \item[Date described] c. 30–60 CE
  \item[Place / stage] Eastern Mediterranean
  \item[Best source or object] \textbf{Paul: minimal anthology}
  \item[Suggested cut] Galatians 2:11–21; 3:6–14; 1 Corinthians 7:29–31; 8:4–6; 15:3–8; Romans 9:1–8; 11:25–32
  \item[What happens to the Bible] Scripture is reread around a crucified and risen messiah; gentiles enter Abraham’s family without becoming ordinary Judeans.
  \item[Approximate date written / made] Authentic letters c. 49–62 CE
\end{description}

\chronologysection{c. 30–90 CE}{Gospel of Thomas}{Egypt or Syria}
\begin{description}
  \item[Date described] c. 30–90 CE
  \item[Place / stage] Egypt or Syria
  \item[Best source or object] \textbf{Gospel of Thomas}
  \item[Suggested cut] Sayings 1–3, 13, 22, 77, 114
  \item[What happens to the Bible] The gospel can exist without narrative, crucifixion story, or resurrection sequence: Jesus becomes a collection of secret sayings.
  \item[Approximate date written / made] Core and final form disputed; commonly first–second century CE
\end{description}

\chronologysection{66–73 CE}{Josephus, Jewish War}{Judea / Rome}
\begin{description}
  \item[Date described] 66–73 CE
  \item[Place / stage] Judea / Rome
  \item[Best source or object] \textbf{Josephus, Jewish War}
  \item[Suggested cut] Preface; 6.249–270 on the temple’s burning
  \item[What happens to the Bible] Biblical historiography becomes Greek-Roman eyewitness history. The second temple falls, forcing another enormous reinterpretation.
  \item[Approximate date written / made] c. 75–79 CE
\end{description}

\chronologysection{70 CE interpreted afterward}{4 Ezra}{Judea and diaspora}
\begin{description}
  \item[Date described] 70 CE interpreted afterward
  \item[Place / stage] Judea and diaspora
  \item[Best source or object] \textbf{4 Ezra}
  \item[Suggested cut] 3:1–36; 7:75–101; 12:10–12
  \item[What happens to the Bible] The destruction of 70 is narrated through Ezra after 586. Theodicy and apocalypse fuse the two temple catastrophes.
  \item[Approximate date written / made] c. 90–100 CE
\end{description}

\chronologysection{70 CE interpreted afterward}{2 Baruch}{Judea and diaspora}
\begin{description}
  \item[Date described] 70 CE interpreted afterward
  \item[Place / stage] Judea and diaspora
  \item[Best source or object] \textbf{2 Baruch}
  \item[Suggested cut] 10:1–19; 31:1–5; 77:12–26
  \item[What happens to the Bible] Baruch mourns Zion and sends revelation to the exiles. The dead Jerusalem again survives by becoming a book.
  \item[Approximate date written / made] c. 90–120 CE
\end{description}

\chronologysection{c. 95 CE}{Revelation}{Roman Asia}
\begin{description}
  \item[Date described] c. 95 CE
  \item[Place / stage] Roman Asia
  \item[Best source or object] \textbf{Revelation}
  \item[Suggested cut] 1:1–8; 5:1–10; 12:1–12; 13:1–10; 21:1–5
  \item[What happens to the Bible] Daniel, Ezekiel, Isaiah, Exodus, and Zechariah are fed into a visionary machine aimed at Rome and a coming New Jerusalem.
  \item[Approximate date written / made] c. 90–96 CE, though debated
\end{description}

\chronologysection{Late first century CE}{Luke–Acts}{Mediterranean world}
\begin{description}
  \item[Date described] Late first century CE
  \item[Place / stage] Mediterranean world
  \item[Best source or object] \textbf{Luke–Acts}
  \item[Suggested cut] Luke 1:1–4; Acts 1:1–8; 15:1–29; 17:16–34; 28:23–31
  \item[What happens to the Bible] The Jesus story becomes world history: Jerusalem generates a Greek-speaking movement that ends preaching in Rome.
  \item[Approximate date written / made] c. 80–100 CE
\end{description}

\chronologysection{c. 50–130 CE}{Dead Sea Scroll biblical manuscripts}{Judean Desert}
\begin{description}
  \item[Date described] c. 50–130 CE
  \item[Place / stage] Judean Desert
  \item[Best source or object] \textbf{Dead Sea Scroll biblical manuscripts}
  \item[Suggested cut] Compare Deuteronomy 32:8; Jeremiah’s lengths; Psalm collections
  \item[What happens to the Bible] There was no single perfectly uniform Hebrew Bible. Masoretic-like, Septuagint-like, Samaritan-like, and independent forms coexist.
  \item[Approximate date written / made] Manuscripts mainly third century BCE–second century CE
\end{description}

\chronologysection{132–135 CE}{Bar Kokhba letters and Babatha archive}{Judean Desert}
\begin{description}
  \item[Date described] 132–135 CE
  \item[Place / stage] Judean Desert
  \item[Best source or object] \textbf{Bar Kokhba letters and Babatha archive}
  \item[Suggested cut] One military order; one legal petition
  \item[What happens to the Bible] Biblical Hebrew, Jewish Aramaic, Greek law, rebellion, property, and ordinary life appear together in real documents.
  \item[Approximate date written / made] 132–135 CE
\end{description}

\chronologysection{c. 170 CE}{Tatian’s Diatessaron}{Syria}
\begin{description}
  \item[Date described] c. 170 CE
  \item[Place / stage] Syria
  \item[Best source or object] \textbf{Tatian’s Diatessaron}
  \item[Suggested cut] Opening harmony combining John 1 and the Synoptic birth narratives
  \item[What happens to the Bible] Four competing gospels are cut into one continuous super-gospel—the Dudetheyreonamystic method applied to Jesus.
  \item[Approximate date written / made] c. 160–175 CE
\end{description}

\chronologysection{c. 180 CE}{Irenaeus, Against Heresies}{Gaul / Rome}
\begin{description}
  \item[Date described] c. 180 CE
  \item[Place / stage] Gaul / Rome
  \item[Best source or object] \textbf{Irenaeus, Against Heresies}
  \item[Suggested cut] 3.11.8 on why there must be exactly four gospels
  \item[What happens to the Bible] The church begins fixing the number and identity of authoritative gospels while rejecting alternative Christian libraries.
  \item[Approximate date written / made] c. 180 CE
\end{description}

\chronologysection{c. 200 CE}{Mishnah}{Palestine}
\begin{description}
  \item[Date described] c. 200 CE
  \item[Place / stage] Palestine
  \item[Best source or object] \textbf{Mishnah}
  \item[Suggested cut] Avot 1:1–2; Yadayim 3:5
  \item[What happens to the Bible] Oral Torah is written beside written Torah. Canon and interpretation become institutions, not merely books.
  \item[Approximate date written / made] c. 200 CE
\end{description}

\chronologysection{Second–fourth centuries CE}{Targums}{Jewish and Christian worlds}
\begin{description}
  \item[Date described] Second–fourth centuries CE
  \item[Place / stage] Jewish and Christian worlds
  \item[Best source or object] \textbf{Targums}
  \item[Suggested cut] Targum Onkelos Genesis 4:8; selected expansion in Targum Jonathan
  \item[What happens to the Bible] Translation into Aramaic becomes interpretation: ambiguities are solved, theology clarified, and narrative gaps quietly filled.
  \item[Approximate date written / made] Traditions early; major forms late antique
\end{description}

\chronologysection{Third–fourth centuries CE}{Origen’s Hexapla}{Caesarea}
\begin{description}
  \item[Date described] Third–fourth centuries CE
  \item[Place / stage] Caesarea
  \item[Best source or object] \textbf{Origen’s Hexapla}
  \item[Suggested cut] Surviving synopsis of Hebrew, transliteration, Aquila, Symmachus, Septuagint, Theodotion
  \item[What happens to the Bible] Textual criticism is born on a monumental scale: six versions aligned to expose disagreement and reconstruct the text.
  \item[Approximate date written / made] c. 240 CE
\end{description}

\chronologysection{Fourth century CE}{Codex Vaticanus and Codex Sinaiticus}{Roman Empire}
\begin{description}
  \item[Date described] Fourth century CE
  \item[Place / stage] Roman Empire
  \item[Best source or object] \textbf{Codex Vaticanus and Codex Sinaiticus}
  \item[Suggested cut] Compare their biblical tables of contents
  \item[What happens to the Bible] The scroll-library becomes a giant bound codex. Christian Bibles contain differing arrangements and edges of canon.
  \item[Approximate date written / made] Fourth century CE
\end{description}

\chronologysection{367 CE}{Athanasius, Festal Letter 39}{Alexandria}
\begin{description}
  \item[Date described] 367 CE
  \item[Place / stage] Alexandria
  \item[Best source or object] \textbf{Athanasius, Festal Letter 39}
  \item[Suggested cut] Canon list
  \item[What happens to the Bible] The first surviving list exactly matching the modern twenty-seven-book New Testament distinguishes canonical books from useful church reading.
  \item[Approximate date written / made] 367 CE
\end{description}

\chronologysection{c. 382–405 CE}{Jerome, Vulgate and prologues}{Bethlehem}
\begin{description}
  \item[Date described] c. 382–405 CE
  \item[Place / stage] Bethlehem
  \item[Best source or object] \textbf{Jerome, Vulgate and prologues}
  \item[Suggested cut] Prologus Galeatus; prefaces to Samuel–Kings and Judith
  \item[What happens to the Bible] Jerome returns to Hebrew, distinguishes canonical from ecclesiastical books, and creates Latin Christianity’s dominant Bible.
  \item[Approximate date written / made] Late fourth–early fifth century CE
\end{description}

\chronologysection{c. 400 CE}{Augustine, Confessions and City of God}{North Africa}
\begin{description}
  \item[Date described] c. 400 CE
  \item[Place / stage] North Africa
  \item[Best source or object] \textbf{Augustine, Confessions and City of God}
  \item[Suggested cut] Confessions 12 on Genesis; City of God 17–18
  \item[What happens to the Bible] The Bible becomes a universal philosophy of history and a text whose obscurities support multiple legitimate readings.
  \item[Approximate date written / made] c. 397–426 CE
\end{description}

\chronologysection{c. 450 CE}{Jerusalem Talmud}{Galilee}
\begin{description}
  \item[Date described] c. 450 CE
  \item[Place / stage] Galilee
  \item[Best source or object] \textbf{Jerusalem Talmud}
  \item[Suggested cut] One brief sugya on textual interpretation
  \item[What happens to the Bible] Scripture is absorbed into legal argument, remembered debate, storytelling, and institutional reasoning.
  \item[Approximate date written / made] Fourth–fifth centuries CE
\end{description}

\chronologysection{c. 500–650 CE}{Babylonian Talmud}{Babylonia}
\begin{description}
  \item[Date described] c. 500–650 CE
  \item[Place / stage] Babylonia
  \item[Best source or object] \textbf{Babylonian Talmud}
  \item[Suggested cut] Bava Batra 14b–15a on authorship and book order; Menachot 29b on Moses and Akiva
  \item[What happens to the Bible] Rabbis explicitly discuss who wrote biblical books and imagine Moses unable to understand later Torah interpretation.
  \item[Approximate date written / made] Fifth–seventh centuries CE
\end{description}

\chronologysection{c. 610–632 CE}{Qur’an}{Arabia}
\begin{description}
  \item[Date described] c. 610–632 CE
  \item[Place / stage] Arabia
  \item[Best source or object] \textbf{Qur’an}
  \item[Suggested cut] Surahs 2:40–79; 5:44–48; 12:1–7; 18:60–82
  \item[What happens to the Bible] Biblical characters and traditions are proclaimed anew in Arabic. The Qur’an confirms earlier revelation while accusing communities of distortion or concealment.
  \item[Approximate date written / made] c. 610–632 CE
\end{description}

\chronologysection{Seventh–tenth centuries CE}{Masoretes}{Tiberias, Babylon, Palestine}
\begin{description}
  \item[Date described] Seventh–tenth centuries CE
  \item[Place / stage] Tiberias, Babylon, Palestine
  \item[Best source or object] \textbf{Masoretes}
  \item[Suggested cut] A page showing vowels, accents, qere/ketiv, marginal Masorah
  \item[What happens to the Bible] Scribes freeze pronunciation, chanting, spelling traditions, and textual statistics around the inherited consonantal text.
  \item[Approximate date written / made] Mainly seventh–tenth centuries CE
\end{description}

\chronologysection{c. 930 CE}{Aleppo Codex}{Tiberias}
\begin{description}
  \item[Date described] c. 930 CE
  \item[Place / stage] Tiberias
  \item[Best source or object] \textbf{Aleppo Codex}
  \item[Suggested cut] One complete annotated page
  \item[What happens to the Bible] The Ben Asher Masoretic system reaches its most prestigious surviving form: consonants, vowels, accents, and scribal guardrails.
  \item[Approximate date written / made] c. 920–930 CE
\end{description}

\chronologysection{1008/1009 CE}{Leningrad Codex}{Cairo}
\begin{description}
  \item[Date described] 1008/1009 CE
  \item[Place / stage] Cairo
  \item[Best source or object] \textbf{Leningrad Codex}
  \item[Suggested cut] Colophon and one representative page
  \item[What happens to the Bible] The oldest complete surviving Hebrew Bible becomes the principal base manuscript of modern critical Hebrew editions.
  \item[Approximate date written / made] 1008/1009 CE
\end{description}

\chronologysection{c. 1080 CE}{Ibn Ezra}{Muslim Spain}
\begin{description}
  \item[Date described] c. 1080 CE
  \item[Place / stage] Muslim Spain
  \item[Best source or object] \textbf{Ibn Ezra}
  \item[Suggested cut] Comments on Deuteronomy 1:1; 34:1–12; “the secret of the twelve”
  \item[What happens to the Bible] A traditional commentator hints that several Torah verses—and perhaps a larger section—were written after Moses.
  \item[Approximate date written / made] Twelfth century CE
\end{description}

\chronologysection{1190 CE}{Maimonides}{Egypt}
\begin{description}
  \item[Date described] 1190 CE
  \item[Place / stage] Egypt
  \item[Best source or object] \textbf{Maimonides}
  \item[Suggested cut] Guide 2.25; Mishneh Torah, Foundations 8
  \item[What happens to the Bible] Philosophical reason constrains interpretation, but Mosaic authorship and Torah’s unique authority remain doctrinal boundaries.
  \item[Approximate date written / made] Late twelfth century CE
\end{description}

\chronologysection{1450s}{Gutenberg Bible}{Mainz}
\begin{description}
  \item[Date described] 1450s
  \item[Place / stage] Mainz
  \item[Best source or object] \textbf{Gutenberg Bible}
  \item[Suggested cut] Opening page and colophon
  \item[What happens to the Bible] Scripture becomes mechanically reproducible. Textual uniformity, mass ownership, vernacular translation, and scholarly comparison accelerate.
  \item[Approximate date written / made] c. 1454–1455
\end{description}

\chronologysection{1516–1517}{First Rabbinic Bible}{Venice}
\begin{description}
  \item[Date described] 1516–1517
  \item[Place / stage] Venice
  \item[Best source or object] \textbf{First Rabbinic Bible}
  \item[Suggested cut] A page with Hebrew text, Targum, and commentaries
  \item[What happens to the Bible] The Jewish Bible appears as a printed conversation among text, Aramaic translation, Masorah, and medieval commentators.
  \item[Approximate date written / made] 1516–1517
\end{description}

\chronologysection{1517–1534}{Martin Luther}{Germany}
\begin{description}
  \item[Date described] 1517–1534
  \item[Place / stage] Germany
  \item[Best source or object] \textbf{Martin Luther}
  \item[Suggested cut] Preface to Romans; prefaces to Hebrews, James, Jude, Revelation; German Bible
  \item[What happens to the Bible] Translation remakes national language; canon develops an internal hierarchy; justification by faith becomes Luther’s interpretive key.
  \item[Approximate date written / made] 1522–1534
\end{description}

\chronologysection{1524–1525}{Second Rabbinic Bible, edited by Jacob ben Hayyim}{Europe}
\begin{description}
  \item[Date described] 1524–1525
  \item[Place / stage] Europe
  \item[Best source or object] \textbf{Second Rabbinic Bible, edited by Jacob ben Hayyim}
  \item[Suggested cut] Masoretic apparatus page
  \item[What happens to the Bible] Printed editorial work standardizes the received Hebrew text for centuries and becomes the base of many later Christian editions.
  \item[Approximate date written / made] 1524–1525
\end{description}

\chronologysection{1546}{Council of Trent canon decree}{Trent}
\begin{description}
  \item[Date described] 1546
  \item[Place / stage] Trent
  \item[Best source or object] \textbf{Council of Trent canon decree}
  \item[Suggested cut] Brief canon list and decree
  \item[What happens to the Bible] In reaction to Protestant exclusions, the Catholic Church dogmatically fixes its larger Old Testament canon.
  \item[Approximate date written / made] 1546
\end{description}

\chronologysection{1651}{Thomas Hobbes, Leviathan}{England}
\begin{description}
  \item[Date described] 1651
  \item[Place / stage] England
  \item[Best source or object] \textbf{Thomas Hobbes, Leviathan}
  \item[Suggested cut] Chapter 33
  \item[What happens to the Bible] Hobbes asks from internal evidence when biblical books could actually have been written, reopening historical authorship in modern Europe.
  \item[Approximate date written / made] 1651
\end{description}

\chronologysection{1670}{Baruch Spinoza, Theological-Political Treatise}{Netherlands}
\begin{description}
  \item[Date described] 1670
  \item[Place / stage] Netherlands
  \item[Best source or object] \textbf{Baruch Spinoza, Theological-Political Treatise}
  \item[Suggested cut] Chapters 7–10, especially discussion of Ezra
  \item[What happens to the Bible] Spinoza treats scripture like any ancient document: Moses did not write the whole Torah; Ezra or a later historian assembled Genesis–Kings.
  \item[Approximate date written / made] 1670
\end{description}

\chronologysection{1678}{Richard Simon, Critical History of the Old Testament}{France}
\begin{description}
  \item[Date described] 1678
  \item[Place / stage] France
  \item[Best source or object] \textbf{Richard Simon, Critical History of the Old Testament}
  \item[Suggested cut] Preface and opening chapters
  \item[What happens to the Bible] Variants, scribes, public recorders, and editorial transmission become central evidence. Catholic textual criticism challenges both Protestant and traditional certainty.
  \item[Approximate date written / made] 1678
\end{description}

\chronologysection{c. 1660–1720}{Isaac Newton’s biblical manuscripts}{England}
\begin{description}
  \item[Date described] c. 1660–1720
  \item[Place / stage] England
  \item[Best source or object] \textbf{Isaac Newton’s biblical manuscripts}
  \item[Suggested cut] \emph{Observations upon the Prophecies} selections; notes on textual corruption and the Trinity
  \item[What happens to the Bible] One of the founders of modern physics treats prophecy, chronology, temple measurements, and textual corruption as an exact historical science.
  \item[Approximate date written / made] Mostly late seventeenth–early eighteenth century; published 1733
\end{description}

\chronologysection{1753}{Jean Astruc, Conjectures}{France}
\begin{description}
  \item[Date described] 1753
  \item[Place / stage] France
  \item[Best source or object] \textbf{Jean Astruc, Conjectures}
  \item[Suggested cut] Opening division of Genesis into sources using divine names
  \item[What happens to the Bible] Repeated stories and alternating divine names are separated into parallel documents. Source criticism receives its basic tool.
  \item[Approximate date written / made] 1753
\end{description}

\chronologysection{1780s–1820s}{Eichhorn and early higher criticism}{Germany}
\begin{description}
  \item[Date described] 1780s–1820s
  \item[Place / stage] Germany
  \item[Best source or object] \textbf{Eichhorn and early higher criticism}
  \item[Suggested cut] Brief methodological preface
  \item[What happens to the Bible] Astruc’s documents cease to be merely Moses’ source notes and become evidence for a long literary history of the Pentateuch.
  \item[Approximate date written / made] Late eighteenth–early nineteenth centuries
\end{description}

\chronologysection{1805}{W. M. L. de Wette}{Germany}
\begin{description}
  \item[Date described] 1805
  \item[Place / stage] Germany
  \item[Best source or object] \textbf{W. M. L. de Wette}
  \item[Suggested cut] Thesis connecting Deuteronomy with Josiah’s reform
  \item[What happens to the Bible] A biblical law book is dated by its fit with a narrated historical event. Deuteronomy becomes a late programmatic work, not Mosaic transcript.
  \item[Approximate date written / made] 1805
\end{description}

\chronologysection{1835}{David Friedrich Strauss, Life of Jesus}{Germany}
\begin{description}
  \item[Date described] 1835
  \item[Place / stage] Germany
  \item[Best source or object] \textbf{David Friedrich Strauss, Life of Jesus}
  \item[Suggested cut] Introduction on myth
  \item[What happens to the Bible] Historical criticism crosses decisively into the Gospels. Supernatural narrative is analyzed as communal myth rather than eyewitness report.
  \item[Approximate date written / made] 1835
\end{description}

\chronologysection{1859–1862}{Darwin and Colenso}{England}
\begin{description}
  \item[Date described] 1859–1862
  \item[Place / stage] England
  \item[Best source or object] \textbf{Darwin and Colenso}
  \item[Suggested cut] Colenso’s arithmetic objections to the Pentateuch
  \item[What happens to the Bible] Numerical impossibilities, historical development, and evolutionary thought make static Mosaic authorship increasingly untenable.
  \item[Approximate date written / made] 1862 onward
\end{description}

\chronologysection{1868}{Mesha Stele discovered}{Dhiban, Moab}
\begin{description}
  \item[Date described] 1868
  \item[Place / stage] Dhiban, Moab
  \item[Best source or object] \textbf{Mesha Stele discovered}
  \item[Suggested cut] Lines on Omri, Israel, Moab, and YHWH
  \item[What happens to the Bible] For the first time, Iron Age neighbors speak in their own voice about the same kingdoms and god found in Kings.
  \item[Approximate date written / made] Inscribed ninth century BCE; discovered 1868
\end{description}

\chronologysection{1866–1878}{Graf, Kuenen, Wellhausen}{Europe}
\begin{description}
  \item[Date described] 1866–1878
  \item[Place / stage] Europe
  \item[Best source or object] \textbf{Graf, Kuenen, Wellhausen}
  \item[Suggested cut] Wellhausen, \emph{Prolegomena}, introduction and temple-history argument
  \item[What happens to the Bible] J and E are early; D belongs near Josiah; P is exilic or postexilic. The Torah becomes the endpoint of Israel’s history rather than its beginning.
  \item[Approximate date written / made] Classical synthesis 1860s–1878
\end{description}

\chronologysection{1870s–1910s}{Gilgamesh, Atrahasis, Enuma Elish recovered and published}{Mesopotamia and Europe}
\begin{description}
  \item[Date described] 1870s–1910s
  \item[Place / stage] Mesopotamia and Europe
  \item[Best source or object] \textbf{Gilgamesh, Atrahasis, Enuma Elish recovered and published}
  \item[Suggested cut] Flood-tablet parallels with Genesis 6–9
  \item[What happens to the Bible] Genesis enters the literature of the ancient Near East. Creation and flood stories are no longer culturally isolated.
  \item[Approximate date written / made] Ancient texts second–first millennia BCE; deciphered chiefly nineteenth century
\end{description}

\chronologysection{1881–1894}{Revised Version and critical Greek New Testament}{Britain}
\begin{description}
  \item[Date described] 1881–1894
  \item[Place / stage] Britain
  \item[Best source or object] \textbf{Revised Version and critical Greek New Testament}
  \item[Suggested cut] Mark 16 and John 7:53–8:11 notes
  \item[What happens to the Bible] Modern translations openly mark passages absent from the earliest manuscripts. The printed Bible begins displaying its own instability.
  \item[Approximate date written / made] Late nineteenth century
\end{description}

\chronologysection{1887–1896}{Amarna letters discovered}{Egypt}
\begin{description}
  \item[Date described] 1887–1896
  \item[Place / stage] Egypt
  \item[Best source or object] \textbf{Amarna letters discovered}
  \item[Suggested cut] Letters from Canaanite rulers asking Pharaoh for help
  \item[What happens to the Bible] Late Bronze Age Canaan speaks directly, complicating biblical conquest stories and illuminating pre-Israelite political geography.
  \item[Approximate date written / made] Fourteenth century BCE; discovered 1887 onward
\end{description}

\chronologysection{1890s–1910s}{Hermann Gunkel and form criticism}{Europe}
\begin{description}
  \item[Date described] 1890s–1910s
  \item[Place / stage] Europe
  \item[Best source or object] \textbf{Hermann Gunkel and form criticism}
  \item[Suggested cut] Introduction to \emph{Genesis}
  \item[What happens to the Bible] Scholars move behind written sources toward oral genres: saga, legend, hymn, law, lament, and the social settings that produced them.
  \item[Approximate date written / made] Major work 1901 onward
\end{description}

\chronologysection{1896–1902}{Herzl and early political Zionism}{Europe / Palestine}
\begin{description}
  \item[Date described] 1896–1902
  \item[Place / stage] Europe / Palestine
  \item[Best source or object] \textbf{Herzl and early political Zionism}
  \item[Suggested cut] \emph{The Jewish State} selections
  \item[What happens to the Bible] The biblical land becomes a modern political destination. Return eventually makes large-scale Hebrew archaeology materially possible.
  \item[Approximate date written / made] 1896
\end{description}

\chronologysection{1920}{Joseph Trumpeldor tradition}{Tel Hai}
\begin{description}
  \item[Date described] 1920
  \item[Place / stage] Tel Hai
  \item[Best source or object] \textbf{Joseph Trumpeldor tradition}
  \item[Suggested cut] Brief account of Tel Hai; cautiously present the disputed final words
  \item[What happens to the Bible] Modern settlement acquires biblical-style martyr legend. Keep this short: it matters mainly as part of the return to the archaeological landscape.
  \item[Approximate date written / made] Event 1920; legend developed afterward
\end{description}

\chronologysection{1928–1934}{Ugaritic tablets}{Syria}
\begin{description}
  \item[Date described] 1928–1934
  \item[Place / stage] Syria
  \item[Best source or object] \textbf{Ugaritic tablets}
  \item[Suggested cut] Baal Cycle selections beside Psalms and Isaiah
  \item[What happens to the Bible] Canaanite poetic grammar, divine councils, storm-god imagery, and parallelism reveal the literary world from which biblical poetry emerged.
  \item[Approximate date written / made] Tablets c. thirteenth century BCE; discovered 1928
\end{description}

\chronologysection{1941–1943}{Martin Noth, Deuteronomistic History}{Europe / United States}
\begin{description}
  \item[Date described] 1941–1943
  \item[Place / stage] Europe / United States
  \item[Best source or object] \textbf{Martin Noth, Deuteronomistic History}
  \item[Suggested cut] Opening methodological statement
  \item[What happens to the Bible] Deuteronomy through Kings is recognized as a coherent edited history explaining national disaster through covenant violation.
  \item[Approximate date written / made] 1943
\end{description}

\chronologysection{1946}{King David Hotel bombing}{Jerusalem}
\begin{description}
  \item[Date described] 1946
  \item[Place / stage] Jerusalem
  \item[Best source or object] \textbf{King David Hotel bombing}
  \item[Suggested cut] A very brief documentary account
  \item[What happens to the Bible] A biblical royal name becomes attached to modern insurgency and empire. Include only as connective tissue, not as the book’s subject.
  \item[Approximate date written / made] Event July 22, 1946
\end{description}

\chronologysection{1947–1956}{Discovery of the Dead Sea Scrolls}{Qumran}
\begin{description}
  \item[Date described] 1947–1956
  \item[Place / stage] Qumran
  \item[Best source or object] \textbf{Discovery of the Dead Sea Scrolls}
  \item[Suggested cut] Isaiah Scroll; Samuel fragments; Jeremiah variants; Deuteronomy 32:8
  \item[What happens to the Bible] Manuscripts a thousand years older than the medieval codices reveal both remarkable continuity and genuine textual plurality.
  \item[Approximate date written / made] Scrolls copied third century BCE–first century CE
\end{description}

\chronologysection{1947–1949}{UN partition, declaration, war, displacement}{Palestine / Israel}
\begin{description}
  \item[Date described] 1947–1949
  \item[Place / stage] Palestine / Israel
  \item[Best source or object] \textbf{UN partition, declaration, war, displacement}
  \item[Suggested cut] Ben-Gurion’s declaration plus one Palestinian and one Jewish eyewitness excerpt
  \item[What happens to the Bible] A Jewish state returns Hebrew institutions to the land, while war creates another mass displacement. Keep this concise and multiperspectival.
  \item[Approximate date written / made] Events 1947–1949
\end{description}

\chronologysection{1948 onward}{David Ben-Gurion and the Bible}{Israel}
\begin{description}
  \item[Date described] 1948 onward
  \item[Place / stage] Israel
  \item[Best source or object] \textbf{David Ben-Gurion and the Bible}
  \item[Suggested cut] Brief remarks from his biblical study circle or speeches
  \item[What happens to the Bible] The Bible is used as national archive, land deed, cultural epic, and stimulus for archaeology—sometimes collapsing ancient and modern Israel into one story.
  \item[Approximate date written / made] Mid-twentieth century
\end{description}

\chronologysection{1948–1957}{W. F. Albright and biblical archaeology}{United States}
\begin{description}
  \item[Date described] 1948–1957
  \item[Place / stage] United States
  \item[Best source or object] \textbf{W. F. Albright and biblical archaeology}
  \item[Suggested cut] One programmatic excerpt
  \item[What happens to the Bible] Archaeology is recruited to verify a broadly historical Bible, although later archaeology becomes less apologetic and more anthropological.
  \item[Approximate date written / made] Mid-twentieth century
\end{description}

\chronologysection{1950s–1970s}{Frank Moore Cross}{United States}
\begin{description}
  \item[Date described] 1950s–1970s
  \item[Place / stage] United States
  \item[Best source or object] \textbf{Frank Moore Cross}
  \item[Suggested cut] \emph{Canaanite Myth and Hebrew Epic} selections
  \item[What happens to the Bible] Source criticism, epigraphy, textual criticism, and Canaanite literature are combined into a history of Israelite religion and biblical tradition.
  \item[Approximate date written / made] Major synthesis 1973
\end{description}

\chronologysection{1967}{Six-Day War}{Jerusalem and surrounding states}
\begin{description}
  \item[Date described] 1967
  \item[Place / stage] Jerusalem and surrounding states
  \item[Best source or object] \textbf{Six-Day War}
  \item[Suggested cut] Brief neutral chronology; one passage on capture of East Jerusalem
  \item[What happens to the Bible] Israeli control expands over many central biblical sites, transforming excavation, tourism, settlement, and the politics of biblical geography.
  \item[Approximate date written / made] June 1967
\end{description}

\chronologysection{1973}{Yom Kippur War}{Israel and surrounding states}
\begin{description}
  \item[Date described] 1973
  \item[Place / stage] Israel and surrounding states
  \item[Best source or object] \textbf{Yom Kippur War}
  \item[Suggested cut] Two or three sentences only
  \item[What happens to the Bible] Biblical festival names continue framing modern events, but modern national history remains secondary to the history of the text.
  \item[Approximate date written / made] October 1973
\end{description}

\chronologysection{1970s–1990s}{Collapse of the old Wellhausen consensus}{Europe and North America}
\begin{description}
  \item[Date described] 1970s–1990s
  \item[Place / stage] Europe and North America
  \item[Best source or object] \textbf{Collapse of the old Wellhausen consensus}
  \item[Suggested cut] Rolf Rendtorff, John Van Seters, Hans Heinrich Schmid
  \item[What happens to the Bible] Many scholars reject four continuously traceable documents, favoring supplements, fragments, blocks, and increasingly late composition.
  \item[Approximate date written / made] 1970s onward
\end{description}

\chronologysection{1981–1987}{Richard Elliott Friedman}{United States / Israel}
\begin{description}
  \item[Date described] 1981–1987
  \item[Place / stage] United States / Israel
  \item[Best source or object] \textbf{Richard Elliott Friedman}
  \item[Suggested cut] \emph{Who Wrote the Bible?} source map and argument for D’s editor
  \item[What happens to the Bible] Classical source criticism is restated as readable detective history, with authors, institutions, politics, and redactors restored to view.
  \item[Approximate date written / made] 1980s
\end{description}

\chronologysection{1993–1994}{Tel Dan Stele discovered}{Tel Dan}
\begin{description}
  \item[Date described] 1993–1994
  \item[Place / stage] Tel Dan
  \item[Best source or object] \textbf{Tel Dan Stele discovered}
  \item[Suggested cut] “House of David” lines
  \item[What happens to the Bible] An Aramaean royal inscription supplies the earliest generally accepted extrabiblical reference to a Davidic dynasty.
  \item[Approximate date written / made] Inscribed ninth century BCE; discovered 1993–1994
\end{description}

\chronologysection{1990s–2010s}{Minimalism, maximalism, and low chronology debates}{Israel / Palestine}
\begin{description}
  \item[Date described] 1990s–2010s
  \item[Place / stage] Israel / Palestine
  \item[Best source or object] \textbf{Minimalism, maximalism, and low chronology debates}
  \item[Suggested cut] Finkelstein, Dever, Thompson, Lemche, Mazar—one paragraph each
  \item[What happens to the Bible] Archaeology stops serving as a simple lie detector. Scholars debate state formation, David and Solomon, ethnic emergence, and the scale of early Israel.
  \item[Approximate date written / made] Late twentieth–early twenty-first centuries
\end{description}

\chronologysection{2000s–present}{Neo-Documentary Hypothesis}{North America}
\begin{description}
  \item[Date described] 2000s–present
  \item[Place / stage] North America
  \item[Best source or object] \textbf{Neo-Documentary Hypothesis}
  \item[Suggested cut] Baruch Schwartz, Joel Baden, Jeffrey Stackert
  \item[What happens to the Bible] A renewed source model identifies documents primarily through incompatible plots and legal systems, not merely vocabulary or divine names.
  \item[Approximate date written / made] Late twentieth–twenty-first centuries
\end{description}

\chronologysection{2000s–present}{Supplementary and redactional models}{Europe}
\begin{description}
  \item[Date described] 2000s–present
  \item[Place / stage] Europe
  \item[Best source or object] \textbf{Supplementary and redactional models}
  \item[Suggested cut] Konrad Schmid, Jan Christian Gertz, Thomas Römer, Reinhard Kratz
  \item[What happens to the Bible] Much European scholarship treats the Pentateuch as successive expansions of major literary cores rather than four complete parallel documents.
  \item[Approximate date written / made] Contemporary scholarship
\end{description}

\chronologysection{2000s–present}{Literary, feminist, postcolonial, archaeological, and social-scientific criticism}{Global scholarship}
\begin{description}
  \item[Date described] 2000s–present
  \item[Place / stage] Global scholarship
  \item[Best source or object] \textbf{Literary, feminist, postcolonial, archaeological, and social-scientific criticism}
  \item[Suggested cut] One exemplary short piece from each method
  \item[What happens to the Bible] The Bible is no longer reconstructed through source criticism alone. Gender, empire, class, trauma, reception, cognition, scribal culture, and material life enter the history.
  \item[Approximate date written / made] Contemporary scholarship
\end{description}

\chronologysection{2010s–present}{Digital manuscripts and computational textual criticism}{Digital world}
\begin{description}
  \item[Date described] 2010s–present
  \item[Place / stage] Digital world
  \item[Best source or object] \textbf{Digital manuscripts and computational textual criticism}
  \item[Suggested cut] Leon Levy Scrolls Library; manuscript collation; multispectral images
  \item[What happens to the Bible] Anyone can compare fragments and codices. Imaging reveals erased or illegible letters; databases make variants searchable at scale.
  \item[Approximate date written / made] Twenty-first century
\end{description}

\chronologysection{Present}{Current Pentateuchal debate}{No single center}
\begin{description}
  \item[Date described] Present
  \item[Place / stage] No single center
  \item[Best source or object] \textbf{Current Pentateuchal debate}
  \item[Suggested cut] A deliberately unresolved final symposium
  \item[What happens to the Bible] Broad agreement: the Torah is composite, internally stratified, and non-Mosaic in its final form. No agreement: exactly how many sources, when they lived, or how editors combined them.
  \item[Approximate date written / made] Ongoing
\end{description}

\chronologysection{Present}{The Dudetheyreonamystic History itself}{The reader’s desk}
\begin{description}
  \item[Date described] Present
  \item[Place / stage] The reader’s desk
  \item[Best source or object] \textbf{The Dudetheyreonamystic History itself}
  \item[Suggested cut] II Kings 25 → scribes → translators → sects → critics → excavators
  \item[What happens to the Bible] The Bible finally becomes the subject of its own biblical history: a library whose contradictions are evidence of the people who preserved rather than erased them.
  \item[Approximate date written / made] Now
\end{description}

\end{document}
